\documentclass[sigconf, 12pt]{acmart}

\usepackage{enumitem}
\usepackage{dirtytalk}
\usepackage{listings}
\usepackage{xcolor}

\title[CS205 Paper 3]{The Ethics of generative AI in commercial software engineering}

\author{Edgar Jose Donoso Mansilla}
\email{e.donosomansilla@digipen.edu}
\affiliation{
	\institution{DigiPen Institute of Technology}
	\country{USA}
}

\begin{document}
\include{src/code-style}

\maketitle

\section{Introduction}
Software engineering is a field that has developed most of today's infrastructure in some way or another, be it via the implementation of a medical device to the websites that allow us to connect the world. However, these projects need to manage several stakeholders, engineers and budgets to make things come to reality. With the improvements in generative AI, there has been a shift towards trying to incorporate this new tool into the software development pipeline. These attempts are not small companies but instead industry leaders such as Microsoft. \cite{microsoft-ai-future}

Given the workings of generative AI, there are several key concerns. The displacement of people for compute centers, the reliability of the software generated and the dilution of understanding the software developed.

\section{Background Information}

\subsection{Core Issue}
This paper will focus on the integration of generative AI in the field of software engineering. There will be 3 key focuses here when it comes down to it's use:
\begin{itemize}
	\item{Displacement of populations}
	\item{Software safety}
	\item{Effects on proprietary and open source software}
\end{itemize}

\subsection{Stakeholders}
\begin{itemize}
	\item{General population}
	\item{Software engineers}
	\item{Companies that use generative AI}
	\item{Companies that make generative AI tools}
\end{itemize}

\section{Current state of generative AI}
This section will focus on the core applications and issues that generative AI is undergoing:

Applications:
\begin{itemize}
	\item{Usage in research}
	\item{Vibe coding}
	\item{Usage in open source projects}
	\item{Data privacy}
\end{itemize}

Issues:
\begin{itemize}
	\item{The ethics of training it on codebases without the creator's explicit approval}
	\item{Hallucinations: The system's inability to tell you that it doesn't know something}
	\item{The "bubble": The financial risks associated with running such an expensive business}
	\item{Population displacement: The space and resources needed to have the compute centers}
	\item{Unreliability: Leaks and the such}
\end{itemize}

\section{Technical Background}
The technical section will focus on the following:
\begin{itemize}
	\item{The math of the models}
	\item{LLMs and how they work}
	\item{Code generation models and what extra features they implement to ensure correct code}
	\item{Why hallucinations occur}
\end{itemize}

\section{Ethical Analysis}

\subsection{Kantianism}
Under Kantianism, the biggest ethical concerns are:
\begin{itemize}
	\item{Using code for training without the owner's conscent (GitHub Copilot)}
	\item{Displacing people to fit more compute centers}
	\item{Using more water to keep compute centers cool}
\end{itemize}

\subsection{Act Utilitarianism}
In act utilitarianism the biggest contrast is whether all the negative effects on other humans helps maximize the happiness in other humans. This can be evaluated through the development of new technologies and research. By evaluating different sources it is possible to evaluate such cases.

\subsection{Virtue Ethics}
Under virtue ethics, the focus is the virtues of the stakeholders and how the decisions they make with this technology and how virtues either align or misalign with them.

\subsection{Social Contract Theory}
This section focuses on what our society expects out of software and how the usage of machine learning can breach the agreements and expectations that we have about the software that we use and how it is made.

\section{Conclusion}

\section{Beyond the surface}

\onecolumn

\bibliographystyle{ACM-Reference-Format}
\bibliography{references}

\end{document}
